\chapter{Choix technologiques et environnement}
\label{ch:techno}

\section{Backend}

Le serveur est construit sur Spring~Boot~3.3.4 et Java~21 (LTS). Cette section justifie ce
choix au regard de l'objectif pédagogique du module, puis détaille les bibliothèques
retenues.

\subsection{Spring Boot comme implémentation moderne des spécifications Jakarta EE}
\label{sec:jee}

Le module dans lequel s'inscrit ce projet porte sur Java Enterprise Edition. Il convient
donc d'expliciter en quoi le recours à Spring~Boot, souvent perçu comme une alternative à
JEE, en constitue en réalité une mise en œuvre.

Java Enterprise Edition --- renommé Jakarta~EE depuis son transfert à la fondation Eclipse
et le passage de l'espace de noms \code{javax} à \code{jakarta} --- n'est pas un produit
mais un ensemble de \emph{spécifications} : JPA pour la persistance, Bean~Validation pour
la validation, l'API~Servlet pour le traitement des requêtes HTTP, CDI pour l'injection de
dépendances, JTA pour les transactions. Un serveur d'applications tel que WildFly ou
GlassFish en fournit une implémentation complète et monolithique. Spring~Boot suit une
autre voie : il implémente ou intègre ces mêmes spécifications, mais les assemble dans un
artefact autonome embarquant son propre conteneur de servlets.

Le projet LIAS s'appuie effectivement sur ces spécifications, et le code le montre de
manière mesurable. Le tableau~\ref{tab:jee} recense, par comptage des importations dans
les sources, l'usage réel de chaque spécification.

\begin{table}[H]
\centering
\caption{Spécifications Jakarta EE effectivement mises en œuvre dans le projet}
\label{tab:jee}
\small
\begin{tabularx}{\textwidth}{@{}p{3.5cm}p{2.1cm}p{3.2cm}X@{}}
\toprule
\textbf{Spécification} & \textbf{Imports} & \textbf{Implémentation} & \textbf{Usage dans LIAS} \\
\midrule
Jakarta Persistence (JPA) & 224 & Hibernate~6.5, via Spring~Data~JPA & \code{@Entity}, \code{@Table}, \code{@Column}, \code{@ManyToOne}, \code{@Enumerated} sur les 31~entités ; JPQL dans les repositories \\
Jakarta Bean Validation & 70 & Hibernate Validator & \code{@NotBlank}, \code{@NotNull}, \code{@Email}, \code{@Size}, \code{@Min} sur les DTO ; \code{@Valid} sur 17~paramètres de contrôleur \\
Jakarta Servlet & 8 & Tomcat~10 embarqué & \code{HttpServletRequest}, \code{FilterChain} dans le filtre d'authentification et le gestionnaire d'erreurs \\
Injection de dépendances & --- & Conteneur Spring (compatible JSR-330) & Injection par constructeur dans l'ensemble des services et contrôleurs \\
Gestion transactionnelle & --- & Spring \code{@Transactional} & Démarcation déclarative, y compris \code{REQUIRES\_NEW} pour l'audit \\
\bottomrule
\end{tabularx}
\end{table}

Trois points méritent d'être développés.

\paragraph{La persistance est de la JPA, pas une abstraction propriétaire.} Les 31 entités
du modèle sont annotées exclusivement avec l'API \code{jakarta.persistence}. Spring~Data
n'ajoute qu'une couche de génération de requêtes au-dessus d'un \code{EntityManager}
standard ; les entités elles-mêmes seraient déployables telles quelles sur un serveur
d'applications JEE. La classe \code{SoftDeleteSafe} intercepte d'ailleurs une exception de
l'API standard, \code{jakarta.persistence.EntityNotFoundException}, pour traiter le cas
d'une association archivée.

\paragraph{Le traitement HTTP repose sur l'API Servlet.} Le filtre
\code{JwtAuthenticationFilter} étend \code{OncePerRequestFilter} et manipule directement
\code{HttpServletRequest} et \code{FilterChain}. Tomcat~10, embarqué par
\code{spring-boot-starter-web}, est un conteneur de servlets conforme à la spécification
Jakarta~Servlet~6. La différence avec un déploiement JEE classique porte sur l'emballage
--- un exécutable autonome plutôt qu'une archive déployée dans un serveur préexistant ---
non sur le modèle de programmation.

\paragraph{La validation est déclarative et standard.} Les contraintes portées par les DTO
relèvent de Bean~Validation. Elles sont évaluées par le conteneur avant l'entrée dans la
méthode de contrôleur, et leurs violations sont converties en réponses~400 structurées par
un gestionnaire d'exceptions centralisé.

Une réserve doit être formulée par honnêteté : la couche web utilise Spring~MVC et non
JAX-RS (Jakarta RESTful Web Services). Les annotations \code{@RestController} et
\code{@GetMapping} sont propres à Spring. C'est la seule spécification JEE significative
qui ne soit pas employée telle quelle, et ce choix est assumé : Spring~MVC s'intègre
nativement avec Spring~Security, dont dépend tout le modèle d'autorisation du projet.

Ce choix architectural sert enfin l'orientation du Master. Une plateforme destinée à
alimenter, à terme, des analyses sur l'activité d'un laboratoire a besoin d'un socle dont
l'écosystème communique avec les outils de traitement de données --- pilotes JDBC
matures, formats d'échange standards, intégration dans des chaînes de traitement. Spring
Boot conserve la rigueur des spécifications Jakarta~EE tout en s'inscrivant dans
l'outillage effectivement utilisé en production, ce qui répond simultanément à l'exigence
académique du module et à la vocation professionnelle de la formation.

\subsection{Bibliothèques retenues}

Le tableau~\ref{tab:stackback} détaille la pile serveur.

\begin{table}[H]
\centering
\caption{Pile technique du backend}
\label{tab:stackback}
\small
\begin{tabularx}{\textwidth}{@{}p{4.3cm}p{1.9cm}X@{}}
\toprule
\textbf{Composant} & \textbf{Version} & \textbf{Rôle dans le projet} \\
\midrule
Java & 21 (LTS) & Langage ; \code{record} utilisé pour tous les DTO \\
Spring Boot & 3.3.4 & Socle applicatif, auto-configuration, serveur embarqué \\
Spring Data JPA / Hibernate & 6.5 & Persistance, \code{@SQLRestriction} pour la suppression logique \\
Spring Security & 6.3 & Filtrage, \code{@EnableMethodSecurity}, \code{@PreAuthorize} \\
PostgreSQL & 18 & SGBD relationnel \\
Flyway & --- & Versionnement du schéma (V1 à V3) \\
jjwt & 0.12.6 & Génération et vérification des jetons JWT (HS512) \\
MapStruct & 1.6.2 & Génération des convertisseurs entité $\rightarrow$ DTO \\
Lombok & --- & Réduction du code répétitif des entités \\
springdoc-openapi & 2.6.0 & Documentation interactive de l'API \\
OpenPDF & 1.3.43 & Génération du rapport d'activité en PDF \\
Spring Boot Starter Mail & --- & Envoi des identifiants et notifications \\
JUnit 5 / Mockito / Spring Security Test & --- & Tests automatisés \\
H2 & --- & Base en mémoire pour les tests d'intégration \\
\bottomrule
\end{tabularx}
\end{table}

\section{Frontend}

L'interface est une application monopage React~19 construite avec Vite~8. Le
tableau~\ref{tab:stackfront} en détaille les composants.

\begin{table}[H]
\centering
\caption{Pile technique du frontend}
\label{tab:stackfront}
\small
\begin{tabularx}{\textwidth}{@{}p{4.3cm}p{1.9cm}X@{}}
\toprule
\textbf{Composant} & \textbf{Version} & \textbf{Rôle dans le projet} \\
\midrule
React & 19.2 & Bibliothèque d'interface, 21 pages et 27 composants \\
Vite & 8.2 & Serveur de développement et outil de construction \\
React Router & 7.18 & Routage client, gardes de route par rôle \\
Axios & 1.19 & Client HTTP, intercepteurs jeton et erreur~401 \\
Tailwind CSS & 4.3 & Système d'utilitaires de style \\
lucide-react & 1.28 & Jeu d'icônes \\
oxlint & 1.75 & Analyse statique \\
\bottomrule
\end{tabularx}
\end{table}

Le document de spécification recommandait par ailleurs TypeScript, Shadcn~UI, Recharts et
Framer~Motion. Ces recommandations n'ont pas été suivies, et il convient de le justifier
plutôt que de le taire. Le tableau~\ref{tab:ecarts} expose les écarts et leurs motifs.

\begin{table}[H]
\centering
\caption{Écarts assumés par rapport aux technologies recommandées}
\label{tab:ecarts}
\small
\begin{tabularx}{\textwidth}{@{}p{2.9cm}p{3.6cm}X@{}}
\toprule
\textbf{Recommandé} & \textbf{Retenu} & \textbf{Justification} \\
\midrule
MySQL & PostgreSQL & Type \code{uuid} natif, contraintes \code{CHECK} expressives, hébergement managé disponible ; le cahier des charges n'impose pas de SGBD \\
TypeScript & JavaScript (ES2022) & Le périmètre du module porte sur le backend ; le contrat d'interface est déjà typé côté serveur par les DTO Java \\
Shadcn UI & Système de composants interne & Sept modules écrits pour le projet ; aucune dépendance externe à maintenir, cohérence visuelle entièrement maîtrisée \\
Recharts & Composant \code{Chart.jsx} interne & Un seul type de graphique est nécessaire ; évite d'embarquer une bibliothèque complète pour un usage marginal \\
Framer Motion & Transitions CSS & Les micro-interactions requises sont couvertes sans dépendance supplémentaire \\
\bottomrule
\end{tabularx}
\end{table}

\section{Base de données}

PostgreSQL a été retenu pour trois raisons techniques : la prise en charge native du type
\code{uuid}, utilisé comme clé primaire de toutes les entités ; l'expressivité des
contraintes \code{CHECK}, qui permettent de garantir au niveau du SGBD la validité des
colonnes portant des énumérations ; et la disponibilité d'une offre hébergée compatible,
qui a permis de travailler sur une base distante sans administration de serveur.

Le schéma n'est jamais généré par l'ORM. Hibernate fonctionne en mode \code{validate} :
il compare les entités au schéma existant et refuse de démarrer en cas d'écart. Les
évolutions passent exclusivement par des migrations Flyway numérotées, appliquées
automatiquement au démarrage. Cette discipline garantit qu'un environnement recréé depuis
zéro obtient exactement le même schéma.

\section{Sécurité}

Les briques de sécurité sont détaillées au chapitre~\ref{ch:securite}. Elles reposent sur
Spring~Security pour le filtrage et l'autorisation, sur la bibliothèque jjwt pour les
jetons signés en HS512, et sur BCrypt avec un facteur de coût de~12 pour le stockage des
mots de passe. Les secrets de connexion ne figurent pas dans le code source : ils sont lus
depuis des variables d'environnement.

\section{Outils de développement}

Le projet est construit avec Maven, dont une distribution est embarquée dans le dépôt afin
que la chaîne de construction ne dépende pas d'une installation préalable. Le contrôle de
version est assuré par Git. La documentation de l'API est générée automatiquement par
springdoc-openapi et consultable via une interface Swagger. Les diagrammes UML sont écrits
en syntaxe Mermaid, ce qui les place sous contrôle de version au même titre que le code et
permet de les rendre en images à la demande --- les figures de ce rapport ont été produites
par ce mécanisme.

\section{Environnement technique}

Le tableau~\ref{tab:env} récapitule l'environnement d'exécution et de validation.

\begin{table}[H]
\centering
\caption{Environnement technique du projet}
\label{tab:env}
\small
\begin{tabularx}{\textwidth}{@{}p{4.3cm}X@{}}
\toprule
\textbf{Élément} & \textbf{Valeur} \\
\midrule
Système d'exploitation & Windows 11 \\
JDK & Eclipse Temurin 21.0.12 (LTS) \\
Gestionnaire de dépendances & Apache Maven 3.9.16 (embarqué dans le dépôt) \\
Exécution Node.js & v24 \\
Base de données & PostgreSQL 18, instance hébergée \\
Port du serveur applicatif & 8080 (configurable) \\
Port du serveur de développement web & 5173 (configurable) \\
Configuration sensible & Variables d'environnement, fichier local exclu du dépôt \\
\bottomrule
\end{tabularx}
\end{table}
